\subsection{Actual coefficient restriction on whole normalized projective spectra}

\begin{lemma}[Homogeneous components commute with actual coefficient maps]
\label{projective:actual-coefficient-homogeneous-components}
\lean{EnvelopeProperty.Berkovich.map_homogeneousComponent}
\leanok
For arbitrary commutative rings $R,S$, every homomorphism $f:R\to S$,
original polynomial $F\in R[z_0,\dots,z_n]$, and $d\in\mathbb N$,
$f_*(F_d)=(f_*F)_d$ for the original total-degree homogeneous components.
\end{lemma}
\begin{proof}
Compare every original monomial coefficient. Both sides equal the image
under $f$ of that coefficient when its total degree is $d$, and zero
otherwise. The zero case uses $f(0)=0$; injectivity is unnecessary.
\end{proof}

\begin{definition}[Whole actual normalized coefficient pullback]
\label{projective:actual-normalized-coefficient-pullback}
\lean{EnvelopeProperty.Berkovich.NormalizedProjectiveSpectrum.coefficientComap, EnvelopeProperty.Berkovich.NormalizedProjectiveSpectrum.coefficientComap_val, EnvelopeProperty.Berkovich.NormalizedProjectiveSpectrum.continuous_coefficientComap}
\leanok
\uses{projective:actual-coefficient-homogeneous-components}
For every such $f$, pullback along $f_*:R[z]\to S[z]$ gives a continuous
map of the WHOLE actual normalized projective seminorm spectra. For every
original $F$, its value is the target seminorm evaluated on $f_*F$.
\end{definition}
\begin{proof}
Compose the entire original valuation with the actual ring homomorphism.
The previous component identity transfers every graded inequality.
The variables are unchanged, so their original finite maximum remains
one. Continuity follows from continuity of actual spectrum pullback and
of the subtype inclusion. No assertion merely on coordinate values is
used to identify entire points.
\end{proof}

\begin{lemma}[Gauss interchange retains every base form]
\label{projective:actual-gauss-base-form-value}
\lean{EnvelopeProperty.Berkovich.Spectrum.polynomialGauss_val_map_C}
\leanok
For any original polynomial seminorm $x$ over $R$, radius $r\ge0$ and
original base form $F$, the actual coefficient-interchanged Gauss point
satisfies $|C_*F|_{x_r}=|F|_x$.
\end{lemma}
\begin{proof}
The actual whole coefficient-swap isomorphism sends $C_*F$ to the constant
polynomial $C(F)$. Prove this whole-ring identity by constants, addition
and multiplication by each original variable. The genuine Gauss value
of a constant polynomial is its original coefficient value.
\end{proof}

\begin{definition}[Original restriction of an arbitrary Gauss-base point]
\label{projective:actual-relative-polynomial-restriction}
\lean{EnvelopeProperty.Berkovich.RelativeNormalizedProjectiveSpectrum.polynomialRestriction, EnvelopeProperty.Berkovich.RelativeNormalizedProjectiveSpectrum.polynomialRestriction_val, EnvelopeProperty.Berkovich.RelativeNormalizedProjectiveSpectrum.polynomialRestriction_polynomialGauss, EnvelopeProperty.Berkovich.RelativeNormalizedProjectiveSpectrum.continuous_polynomialRestriction}
\leanok
\uses{projective:actual-normalized-coefficient-pullback, projective:actual-gauss-base-form-value}
Let $v$ be any real valuation on $R$ and $r\ge0$. Every actual relative
normalized projective point over the prescribed Gauss valuation $v_r$ on
$R[t]$ restricts continuously to an actual relative normalized projective
point over $v$ on $R$. EVERY original base form has its whole pulled-back
value. Restriction of the actual Gauss extension of an original relative
point recovers the SAME entire original point. The arbitrary input point
is not required to be a Gauss extension.
\end{definition}
\begin{proof}
Apply whole normalized coefficient pullback to $R\to R[t]$. The prescribed
coefficient condition restricts correctly because $v_r(C(a))=v(a)$.
For the recovery statement, compare every original form and apply the
actual Gauss base-form identity, then extensionality of whole valuations
and the two actual subtype objects. Continuity restricts the preceding
whole continuous pullback to the genuine prescribed-coefficient loci.
\end{proof}
